\documentclass[10pt,a4paper]{amsart}
\usepackage[margin=19mm]{geometry}
\usepackage{amsmath,amssymb,mathtools}
\usepackage[scr=rsfs,cal=euler]{mathalfa}
\usepackage{xfrac}
\usepackage[unicode,hidelinks,bookmarks=false]{hyperref}
\newcommand{\ie}{i.e.\ }
\newcommand{\cf}{cf.\ }
\newcommand{\resp}{respectively\ }
\newcommand{\Subsection}{\subsection}
\providecommand{\nobreakdash}{\nobreak\hbox{-}}
\setlength{\emergencystretch}{2em}
\multlinegap0pt
\usepackage{bm}
\usepackage{bbm}
\usepackage{dsfont}
\usepackage{xspace}
\usepackage{cancel}
\usepackage{mathtools}
\usepackage{cleveref}
\usepackage{amsthm}
\makeatletter
\@ifundefined{lemma}{\newtheorem{lemma}{Lemma}}{}
\makeatother
\title[Power maps on General Linear groups over finite principal id]{Power maps on General Linear groups over finite principal ideal local rings of length two}
\author{Saikat Panja, Ayon Roy, Anupam Singh}
\date{Source: arXiv:2508.18952v1 (CC BY 4.0)}
\begin{document}
\maketitle

\noindent\emph{Source and licence.} Power maps on General Linear groups over finite principal ideal local rings of length two. Version arXiv:2508.18952v1 (CC BY 4.0), distributed under the Creative Commons Attribution 4.0 International licence, as recorded in the arXiv licence metadata for this exact version and shown on the abstract page https://arxiv.org/abs/2508.18952v1. Full rights notice: licenses/arxiv-2508.18952-CC-BY-4.0.txt.\par\smallskip

\noindent\textit{Editorial scope.} Counted unit: the Lemma whose source label is \texttt{lem:tech-1-thm-1} in the article (source0.tex lines 250--252, source label \texttt{lem:tech-1-thm-1}), delivered together with the article's own complete original proof of it (source0.tex lines 253--279, one \texttt{proof} environment of 27 lines). No mathematics is written, repaired or re-derived: statement and proof are the article's own text line for line. Required context retained uncounted: lem:snap-factor (lines 195--201); lem:primary-modulo (lines 203--205); hensel-1 (lines 213--215); lem:f-irred-F-irred (lines 233--235). Every definition, display and label of those lines is retained unchanged. Omitted from this package and not redistributed: 1--194, 202--202, 206--212, 216--232, 236--249, 280--1356. Nothing inside a retained excerpt is omitted or altered: every retained body is delivered exactly as the article's lines are written, so no omission receipt of any kind accompanies this package. Citations: the retained excerpts carry no citation command at all, so no bibliography entry of the article is transcribed and no reference is redistributed. Reference adaptation: none was needed and none was made; every reference inside the delivered unit resolves inside it, so no unresolved reference or citation is produced. Printed slips kept literal: no printed word, symbol, hypothesis or number of the counted statement or of its proof was altered. Layout and class: the article's own class is \texttt{amsart} with packages including mathtools, amsmath, amssymb, mathrsfs, geometry, graphicx, hyperref, cleveref, biblatex, lipsum, adjustbox, tikz-cd, mathptmx, float; the wrapper is the lane's pinned \texttt{amsart} editorial wrapper at 10pt on A4, which restates the theorem-like environments the retained text needs and adds the article's own macro definitions that the retained text needs (none), with extra wrapper packages (bm, bbm, dsfont, xspace, cancel, mathtools, cleveref). The delivered numbering is the unit's own. Assets: no retained span contains a figure environment, a graphics-inclusion command, a PDF-page inclusion, a table or a TikZ picture; no image file is copied into this package, the package declares no asset dependency and no image file is redistributed with it. Licence: version arXiv:2508.18952v1 of this article is distributed under the Creative Commons Attribution 4.0 International licence, as recorded in the arXiv licence metadata for this exact version and shown on the abstract page https://arxiv.org/abs/2508.18952v1. A full rights notice is delivered at licenses/arxiv-2508.18952-CC-BY-4.0.txt. Characters: every retained excerpt is pure ASCII, so the delivered engine, XeLaTeX with its default Latin Modern text font, reports no missing character.
\begin{lemma}\label{lem:snap-factor}
Let $\mathscr R$ be a ring, $N(\mathscr R)$ be its nilradical, and $\mathscr R/N(\mathscr R)$ be an integral, principal ideal ring. Then the following are true.
\begin{enumerate}
\item Every non-nilpotent element $\alpha$ of $\mathscr R$ can be factored as $\alpha = \delta \sigma_1\sigma_2\ldots\sigma_n$, where $\delta$ is a unit and $\sigma_1,\ldots,\sigma_n$ are primary, not nilpotent non-units, which are coprime in pairs.
\item If $\delta \sigma_1 \ldots \sigma_n = \delta' \sigma_1' \ldots \sigma_{n'}'$, where $\delta$ and $\delta'$ are units, $\sigma_1, \ldots, \sigma_n$ are primary non-units which are coprime in pairs and the same is true for $\sigma'_1,\ldots, \sigma'_{n'}$; then $n = n'$ and after a suitable reordering, $\sigma_i$ is associated with $\sigma_i'$ for $i=1, 2, \ldots, n$. If $n> 1$, the elements $\sigma_1, \ldots,\sigma_n, \sigma'_1, \ldots, \sigma'_{n'}$ are necessarily not nilpotent.
\end{enumerate}
\end{lemma}

\begin{lemma}\label{lem:primary-modulo}
An element $\alpha\in \mathscr R$ is primary and non-nilpotent if and only if $\overline{\alpha}$ is a primary nonzero element of $\mathscr R/N(\mathscr R)$.
\end{lemma}

\begin{lemma}[Hensel's lemma version 1]\label{hensel-1}
Let $\mathscr R$ be a Noetherian ring, complete with respect to an ideal $\mathfrak{m}$. Let $F(t)$ be a polynomial in $\mathscr R[t]$ and $f(t)$ be the polynomial over $\mathscr R/\mathfrak{m}$ obtained by reducing $F(t)$ modulo $\mathfrak{m}$. Suppose $f(t)$ has a factorization $f(t) = g_1(t)g_2(t)$ in $\mathscr R/\mathfrak{m}[t]$ in such a way that $g_1(t)$ and $g_2(t)$ generate the unit ideal, and $g_1(t)$ is monic. Then, there is a unique factorization $$F(t)=G_1(t)G_2(t)\in \mathscr R[t]$$ such that $G_1(t)$ is monic and $G_i(t)$ reduces to $g_i(t)$ mod $\mathfrak{m}$ for $i=1, 2$.
\end{lemma}

\begin{lemma}\label{lem:f-irred-F-irred}
Let $\mathscr O_\ell$ be a local principal ideal ring of length $\ell$, $\mathfrak{m}= \langle a\rangle$ be its unique maximal ideal, and $k$ be the residue field. Let $F(t)\in \mathscr O_\ell[t]$ be a fundamental irreducible polynomial. Then $F(t)$ must be irreducible in $\mathscr O_\ell[t]$.
\end{lemma}

\begin{lemma}[Unique factorization using  monic fundamental irreducible polynomials]\label{lem:tech-1-thm-1}
Let $\mathscr O_{\ell}$ be a local principal ideal ring of length two and $F(t) \in\mathscr O_\ell[t]$ be a monic polynomial such that $\overline{F}(t) = f(t)$ is a separable polynomial in $k[t]$. Then $F(t)$ has a unique factorization into coprime fundamental irreducibles in $\mathscr O_\ell[t]$, up to a permutation of fundamental irreducible factors.
\end{lemma}

\begin{proof}
We will first show the existence of such a factorization and use \Cref{lem:snap-factor} to prove its uniqueness.

Let $F(t)\in \mathscr O_\ell[t]$ be a monic polynomial such that $\overline{F}(t) = f_1(t)f_2(t)\cdots f_r(t)$, where $f_i(t)$ are distinct irreducibles in $k[t]$ for $1\leq i\leq r$.
Without loss of generality, let us take $f_i(t)$ to be a monic polynomial (if not monic, divide by the leading coefficient) for some fixed $i$ and $g_1(t) = f_2(t)\cdots f_r(t)$.
Clearly $\gcd(f_1(t), g_1(t)) = 1$. 
Then by Hensel lemma version 1, \Cref{hensel-1}, we get polynomials $F_1(t)$ and $G_1(t)$ in $\mathscr O_\ell[t]$ such that $F(t) = F_1(t)G_1(t)$ where $\overline{F}_1(t) =f_1(t)$ and $\overline{G}_1(t) =g_1(t)$ in $k[t]$ and $F_1(t)$ is monic.
Now as $F_1(t)$ is monic and $f_1(t)$ is irreducible in $k[t]$, by \Cref{lem:f-irred-F-irred}, $F_1(t)$ must be irreducible in $\mathscr O_\ell[t]$ and hence is a fundamental irreducible.

Next, consider the polynomial ${g}_2(t) = f_3(t)\cdots f_r(t).$ 
Then, $\gcd(f_2(t), {g}_2(t))=1$. Moreover, $f_2(t)$ is monic irreducible and $g_1(t) = f_2(t){g}_2(t)$. Then by Hensel's lemma version 1, \Cref{hensel-1}, we get a factorization of $G_1(t)$ in $R[t]$ as $G_1(t) = F_2(t)G_2(t)$; where $F_2(t)$ is monic and $\overline{F}_2 =f_2, \overline{G}_2={g_2}$. 
Now, $F_2(t)$ must be fundamental irreducible. Hence, we get $F(t)=F_1(t) F_2(t) G_2(t)$. Continuing this process we achieve a factorization (with ordering) $F(t) =F_1(t)\cdots F_{r-1}(t)F_r(t)$ where each $F_i(t)$ is fundamentally irreducible and $F_1(t),\ldots, F_{r-1}(t)$ are monic. This implies $\deg(F_i(t)) = \deg(f_i(t))$ for $1\leq i\leq r$. As $F(t)$ is monic, this forces $F_r(t)$ to be monic (compare the leading term coefficient, as in \Cref{lem:f-irred-F-irred}). Note that all $F_i(t)$ are pairwise distinct, because otherwise $f_i(t)$ will not be distinct for $1\leq i\leq r$.

Now to prove uniqueness, if possible let there be an another decomposition of $F(t)$ in monic fundamental irreducibles as follows $$F(t)= H_1(t)H_2(t)\cdots H_{r'}(t)$$ such that all $H_i(t)$ are non-constant monic fundamental irreducible in $R[t]$ (this forces $r=r'$) and $\theta({H_i})=f_i$, $\deg(H_i(t))=\deg(f_i(t))$ for each $1\leq i\leq r$.
Consider $A=R[t]$ and $N(A)=\mathfrak{m}[t]$ in \Cref{lem:snap-factor}.
Clearly $F_1(t), F_2(t),\ldots, F_r(t)$ are primary (since $f_i$ is irreducible, hence $\langle f_i\rangle$ is prime, and thus \Cref{lem:primary-modulo} is applicable) non-units and they are relatively divisorless in pairs.
Moreover the same is true for $H_1(t),H_2(t),\ldots,H_r(t)$ also.
Then by \Cref{lem:snap-factor}  we get , after a suitable ordering $\langle F_i(t)\rangle=\langle H_i(t)\rangle$ for $1\leq i\leq r$. 
Therefore $H_i(t)=\alpha_i(t)F_i(t)$ where $\alpha_i(t)\in R[t]$ for $1\leq i\leq r$. 
As $H_i(t)$ is a monic irreducible polynomial, $\alpha_i(t)$ must be a unit in $\mathscr O_\ell[t]$.
This implies that $\alpha_i(t)=a_i+m_i(t)$ for each fixed $i$; where $a_i\in R^{\times}$ and $m_i[t]\in \mathfrak{m}[t].$ 

If possible, let $m_i(t)$ be non-constant.
As $F_i(t)$ is monic therefore the degree of $\alpha_i(t)F_i(t)$ exceeds the degree of $F_i(t)$. Which is a contradiction (since $\deg(H_i(t))$ is same as $\deg(F_i(t))$ by assumption). 
Hence, each $\alpha_i(t)$ should be a constant and that should be $1$.
Hence after the suitable ordering we get $F_i(t)=H_i(t)$ for $1\leq i\leq r$. So the decomposition is unique.
\end{proof}

\end{document}
